\documentclass[11pt]{article}

\usepackage[margin=1in]{geometry}
\usepackage{amsmath,amssymb}
\usepackage{booktabs}
\usepackage{graphicx}
\usepackage{url}
\usepackage{hyperref}
\usepackage{xcolor}

\newcommand{\strata}{STRATA}
\newcommand{\code}[1]{\texttt{#1}}

\title{\strata: Catastrophic Forgetting as an Architectural Default,\\
a Measured Failure Taxonomy, and a Replay-Free Resolution\\
at Prototype Scale}

\author{Christopher~Gardner\thanks{Correspondence:
christhechef35@gmail.com. Architecture design, implementation and
experiments were carried out in an interactive research session with
Claude (Anthropic); all results reported here are reproducible from the
accompanying code.}}

\date{July 2026}

\begin{document}
\maketitle

\begin{abstract}
Catastrophic forgetting is conventionally treated as a training pathology to
be suppressed with regularizers or replay. We argue, and demonstrate at
prototype scale, that it is instead the \emph{default consequence} of three
architectural decisions most networks make implicitly: writing new
information into the same substrate as old information, at the same
timescale, in a correlated basis. We present \strata\ (Sparse Tiered Routing
with Anchored Two-timescale Adaptation), a continual-learning architecture
that makes all three decisions differently, using only local learning rules:
no backpropagation through time, no task labels, and no replay buffers of
raw data. Seven composable mechanisms---sparse hysteretic routing,
atomic feature allocation, a fast/slow synaptic split, multi-timescale
cascade consolidation, maturity-gated readouts, predictive-residual coding,
and inference-side decorrelation with impermanent consolidation---are each
pinned to a specific, empirically measured failure mode. On sequential
task benches with forced cross-task interference, \strata\ exhibits flat
retention over 24--48 sequential tasks, capacity growth that tracks task
\emph{diversity} rather than task count (16 columns for 48 tasks across 16
families), up to $2.6\times$ forward transfer, and graceful power-law
degradation; with representation learning enabled, the learned dictionary
ultimately \emph{outperforms} a fixed random expansion on the primary
bench.
We further embed the architecture as a metabolically constrained organ in an
artificial organism, where learning generates heat that throttles its own
plasticity, and demonstrate a two-layer stack in which event-driven ascent
and bounded top-down priors reduce regime re-orientation transients by
15--31\%; the same experiments yield a measured design constraint for
hierarchy---abstraction density at level $N{+}1$ requires geometrically
more experience than level $N$ can supply per unit time. Alongside the
positive results we report a taxonomy of failure modes---absorption, code
drift, correlation inheritance, capacity exhaustion, orphaned bindings,
threshold collapse, and prior capture---each discovered by measurement,
several of which recurred at higher levels of the hierarchy in
level-specific form, and a falsification record in which every mechanism
was tested against a pre-registered failure criterion. We state the scope
limits plainly: all results are at toy scale on self-designed benches, and
composition under depth beyond two layers, external benchmarks, and
order-of-magnitude scaling remain open.
\end{abstract}

\section{Introduction}

This paper's central claim is that ``catastrophic forgetting'' is not one
phenomenon but several: an umbrella term for a set of distinct,
individually measurable failure modes---\emph{absorption} of an old
module's identity by new data, \emph{code drift} invalidating downstream
bindings, \emph{correlation inheritance} in learned features,
\emph{capacity exhaustion} of finite dictionaries, and \emph{orphaned
episodic bindings} under representation change---each of which can be
detected by instrumentation, prevented by a specific mechanism, and
unit-tested. We present this as a \emph{candidate} decomposition: it is
internally predictive within the systems studied here (\S\ref{sec:falsification}),
and whether it is the structure of the phenomenon at large---or the first
chart of it---is precisely what evaluation on benchmarks we did not design
must decide. What the decomposition does establish is that, in the systems
studied, forgetting is the default consequence of three architectural
decisions most networks make implicitly: writing new information (i) into
the same physical substrate as old information, (ii) at a single timescale
of plasticity, and (iii) in a representational basis whose components are
correlated across experiences.

Modern neural networks learn well from stationary, shuffled data and forget
catastrophically from sequential, non-stationary data
\cite{french1999,mccloskey1989}. The dominant remedies---elastic
regularization of important weights \cite{kirkpatrick2017,zenke2017},
replay of stored or generated past data \cite{shin2017}, or frozen
task-specific expansion \cite{rusu2016}---treat forgetting as a defect of
the \emph{training procedure}. Under the decomposition above they are
partial by construction: a loss term attacks only the timescale condition,
and replay converts the non-stationary problem back to a stationary one at
the cost of storing data, while the interference both penalize is
manufactured upstream of the loss.

We test the hypothesis constructively. \strata\ is a small,
fully-instrumented continual-learning system in which each of the three
decisions above is made differently, by mechanisms with independent support
in theoretical neuroscience: complementary learning systems
\cite{mcclelland1995}, cascade models of synaptic consolidation
\cite{fusi2005,benna2016}, sparse expansion coding \cite{olshausen1996,
dasgupta2017}, predictive coding \cite{rao1999}, and reservoir computing
\cite{maass2002,jaeger2001}. The contribution is not any single mechanism
---most have antecedents---but (a)~evidence that they \emph{compose} into
one working system under strict constraints (local rules only, no replay,
no task labels, bounded state), and (b)~a \emph{measured failure taxonomy}:
a catalogue of the specific ways such systems break, each observed in
instrumented experiments, each answered by a named mechanism, and several
observed to recur at higher levels of a hierarchy in new costume. We
believe the taxonomy is the most transferable artifact, because these
failures will confront any architecture attacking this problem at any
scale.

\paragraph{Contributions.}
\begin{enumerate}
\item An architectural account of catastrophic forgetting (substrate,
timescale, basis) and a seven-mechanism blueprint that operationalizes it
with purely local learning rules (\S\ref{sec:architecture}).
\item A measured failure taxonomy for replay-free continual learning
(\S\ref{sec:taxonomy}), including three failure modes we have not seen
named elsewhere: \emph{absorption} (identity capture without weight
erasure), \emph{correlation inheritance} in learned dictionaries, and
\emph{orphaned episodic bindings} under representation drift.
\item Benchmarks with \emph{forced} interference (shared vocabulary,
conflicting successors), motivated by the measured observation that
disjoint-task benches are jointly fittable and produce no forgetting even
in unprotected baselines (\S\ref{sec:benches}).
\item Empirical results: flat retention across 24 sequential tasks;
capacity growth tracking diversity, not task count; $2.6\times$ forward
transfer; and a learned dictionary that ultimately beats a fixed random
expansion under continual-learning constraints (\S\ref{sec:results}).
\item A metabolic embedding in which synaptic work is paid for in a
conserved energy economy, producing an emergent learn/overheat/rest cadence
and reproducing a known saturation pathology of the host system
(\S\ref{sec:thermo}).
\item A two-layer stack with event-driven ascent and bounded top-down
priors, reducing regime re-orientation transients by 15--31\%, together
with negative results delineating when hierarchy does \emph{not} pay
(\S\ref{sec:depth}).
\end{enumerate}

\section{Forgetting as an Architectural Default}
\label{sec:principles}

Consider online learning of a predictive map from a non-stationary stream.
Interference between a new experience and consolidated knowledge requires
three simultaneous conditions: the update must touch \emph{the same
parameters} (substrate overlap), it must move them \emph{faster than they
can be protected} (timescale collision), and the representations involved
must \emph{overlap} (basis correlation). Standard architectures satisfy all
three by default: dense representations make every parameter serve every
input; a single learning rate makes protection and adaptation the same
knob; and learned features inherit whatever correlations exist in the
input distribution. Regularization methods attack the middle condition
only, and replay methods convert the non-stationary problem back to a
stationary one at the cost of storing data. \strata\ instead removes each
condition architecturally: routing and sparsity decide \emph{where}
information is written; a fast/slow split and consolidation cascades decide
\emph{when}; and predictive-residual coding decides \emph{what}---only the
deviation from a slow expectation ever reaches plastic weights.

\section{The \strata\ Architecture}
\label{sec:architecture}

\strata\ is an online next-input predictor. Inputs $x_t \in \mathbb{R}^d$
arrive one per tick with no task labels and no episode boundaries. All
learning rules are local (delta rules and imprinting); the only
temporal credit assignment is through bounded liquid state, never through
backpropagation through time.

\subsection{Routing and the slow pathway}
Each \emph{column} $c$ owns a centroid $\mu_c \in \mathbb{R}^d$: an
exponential moving average of the inputs it has claimed, serving
simultaneously as (i) its routing identity and (ii) its slow
\emph{expectation} of the input. A slow context $\kappa_t = (1-\rho)\,
\kappa_{t-1} + \rho\, x_t$ is routed by cosine affinity
$a_c = \cos(\kappa_t, \mu_c)$ with \textbf{three-way hysteresis}:
\begin{align}
a_c \ge \theta_{\mathrm{claim}} &: \text{column may absorb the input into
its identity (and update claim statistics)}\nonumber\\
a_c \ge \theta_{\mathrm{route}} &: \text{column participates in prediction
\emph{and} learning}\nonumber\\
a_c < \theta_{\mathrm{route}} &: \text{column may serve as fallback
predictor but receives \emph{zero} learning,}\nonumber
\end{align}
with $\theta_{\mathrm{claim}} \gg \theta_{\mathrm{route}}$ and all identity
updates frozen while a regime transition is detected. Persistent orphanhood
(no column routable), measured by a leaky accumulator rather than a
consecutive-tick counter, spawns a new column warm-started from its nearest
neighbour. Optionally, thresholds are \emph{adaptive}: per-column claim
statistics $(m_c, v_c)$ tracked only on claimed ticks give
$\theta_{\mathrm{claim},c} = \max\!\big(m_c - \lambda\sqrt{v_c},\;
b + \kappa\,\sigma_b\big)$, where $(b,\sigma_b)$ are background-similarity
statistics tracked on orphaned ticks; the floor term is load-bearing
(\S\ref{sec:taxonomy}, \emph{threshold collapse}).

\subsection{Predictive-residual coding (the two-pathway split)}
The plastic pathway never sees the raw input. With $\mu$ the claiming
column's centroid, the residual $r_t = x_t - \mu$ is encoded by a sparse
expansion $z_t = \mathrm{enc}(r_t) \in \mathbb{R}^{D}$ ($s$-sparse,
$s \ll D$), and the final prediction adds the expectation back:
\begin{equation}
\hat{x}_{t+1} \;=\; \mu \;+\; \underbrace{\textstyle\sum_{c} g_c \big(
C_c h^c_t + D_c z_t + Q_c q_t \big)}_{\text{slow readouts}}
\;\oplus\; \underbrace{F z_t}_{\text{fast weights}},
\end{equation}
where $\oplus$ denotes confidence-weighted blending. Because the shared
(correlated) component of a regime is absorbed by $\mu$, the dictionary
codes only deviations---removing, by construction, the input correlations
that learned features were measured to inherit fatally
(\S\ref{sec:encoderarc}).

\subsection{Liquid state and readouts}
Each column carries a fixed, diagonal, input-dependent leaky reservoir with
\emph{liquid} time constants stretched by novelty $\nu_t$:
\begin{equation}
h^c_t = e^{-1/\tau_c(x_t,\nu_t)} \odot h^c_{t-1} +
\big(1 - e^{-1/\tau_c(x_t,\nu_t)}\big) \odot (W_B x_t),
\qquad \tau_c = \tau_{\min} + \mathrm{softplus}(w_\tau^{\!\top} x_t)
\,(1 + \alpha\,\nu_t).
\end{equation}
The decay factor lies in $(0,1)$ elementwise, so the state is bounded for
all time. Three linear readouts learn by local delta rules: $C_c$ over the
state, $D_c$ over the code, and $Q_c$ over a fixed random $k$WTA
\emph{conjunctive expansion} $q_t = \mathrm{kWTA}_S\!\big(R\,[z_t \oplus
h^c_t]\big)$, which manufactures symbol-in-temporal-context features and
lifts effective key rank above low-rank input manifolds at zero churn risk
(nothing in $R$ ever moves).

\subsection{The fast/slow split and consolidation}
A fast-weight matrix $F$ binds keys to residual targets by a
novelty-gated delta rule, $F \leftarrow F + \eta^w_t\,(v_t - F z_t)
z_t^{\!\top}$, capturing surprising experience one-shot with no gradient.
During calm periods, anchors sampled from $F$'s own content are distilled
into the claiming column's slow readout: consolidation happens later, from
the system's own memory, never from stored raw data. Every slow parameter
is the surface of a Benna--Fusi diffusion chain $u_1,\dots,u_m$ with
geometrically spaced timescales; learning writes only $u_1$, and symmetric
inter-level exchange converts exponential forgetting into power-law
retention with relearning savings \cite{benna2016,fusi2005}.

\subsection{Learning the dictionary}
The encoder's rows may themselves learn, under three protections:
\textbf{(i) per-unit consolidation}, $\eta_i = \eta_0 / (1 + \kappa\,
\mathrm{usage}_i)$; \textbf{(ii) novelty recruitment with atomic
allocation}: when no row matches a residual, the least-used row is drafted
and \emph{imprinted} one-shot on the Gram--Schmidt residual of the input
against co-active mature rows, and the input is \emph{re-encoded before
anything downstream binds to it}, so keys are born stable; and
\textbf{(iii) usage decay}, which makes consolidation impermanent so a
finite dictionary forgets gracefully, least-used first. Inference may run
either as $k$WTA or as \emph{matching pursuit} \cite{mallat1993}: each
selected atom's contribution is subtracted before the next selection, which
performs lateral decorrelation \emph{by inference} rather than by weight
dynamics and therefore cannot conflict with consolidation. Slow readouts
consolidate onto a key only as fast as the key's supporting rows stop
moving (\emph{maturity gating}). Dictionaries may also grow structurally:
capacity added under allocation pressure is \emph{dormant}---masked out of
inference entirely until the draft-and-imprint path activates it---which
makes growth exactly backward-transparent (bit-identical old codes,
unit-tested).

\subsection{The master signal}
A single standardized surprisal $\nu_t = \|x_t - \hat{x}_t\|^2 /
\widehat{\mathrm{Var}}$ gates everything: fast-weight writes open with
novelty; slow learning closes with it ($1/(1+\max(0,\nu_t-1))$), so
surprising samples are captured episodically rather than consolidated;
a CUSUM statistic on $\nu_t$ declares regime transitions, freezing identity
updates and damping mature columns; and sustained routing orphanhood drives
structural growth.

\section{A Measured Failure Taxonomy}
\label{sec:taxonomy}

Every entry below was observed in an instrumented experiment before being
named; each is prevented by a specific mechanism and, where marked
($\uparrow$), was observed to \emph{recur at the next level of the
hierarchy} in level-specific form.

\begin{table}[ht]
\centering\small
\begin{tabular}{p{3.2cm} p{5.6cm} p{5.6cm}}
\toprule
\textbf{Failure} & \textbf{Signature (measured)} & \textbf{Countermeasure}\\
\midrule
Absorption $\uparrow$ & A new task drags an old column's centroid onto
itself; affinity never dips below threshold (0.62 minimum at a boundary;
0.92 across a schedule change at layer 2) & Claim $\gg$ route hysteresis;
identity frozen during transitions; fallback predicts but never learns\\
\addlinespace
Code drift & Learned encoder rotation invalidates downstream keys;
old-task error $0.23 \to 0.76$ & Per-unit consolidation; novelty
recruitment; two-pathway residual coding\\
\addlinespace
Correlation inheritance & Learned features fire for every correlated
sibling; within-task code separation $|\!\cos\!|$ $0.21 \to 0.33$--$0.45$
& Residual coding (shared structure never reaches the dictionary);
matching-pursuit inference\\
\addlinespace
Orphaned bindings $\uparrow$ & Episodic keys formed pre-allocation never
match post-imprint codes & Atomic allocation: re-encode after imprint;
layer boundaries transmit only post-settlement codes\\
\addlinespace
Capacity exhaustion & All virgin rows consumed (288 symbols vs.\ 128
units); floors degrade $0.168 \to 0.258$ while a random encoder stays flat
& Usage decay (graceful, least-used-first); dormant structural growth\\
\addlinespace
Threshold collapse & A column claiming noise widens its own adaptive band;
absorption returns through the statistics & Background floor
$\theta_{\mathrm{claim}} \ge b + \kappa\sigma_b$; statistics update only on
claimed ticks\\
\addlinespace
Prior capture & A confident-but-wrong top-down prior forces the lower layer
into the wrong regime & Bounded prior bonus (can promote a routable column,
never make an alien one routable); priors bias inference, never learning\\
\addlinespace
Consolidation mis-attribution & Key-affinity anchor assignment distilled
one task's memories into another task's column & Ownership from the routing
context, not key similarity\\
\bottomrule
\end{tabular}
\caption{The failure taxonomy. Numbers reference the experiments of
\S\ref{sec:results}.}
\label{tab:taxonomy}
\end{table}

\section{Benchmarks}
\label{sec:benches}

All benches stream noisy cyclic sequences of unit vectors
($d{=}24$, cycle lengths 12--13, additive noise $\sigma{=}0.02$), one
sample per tick, with tasks switched without any signal to the learner.
A methodological finding shaped their design: \textbf{disjoint-vocabulary
task sequences produce no measurable forgetting even in unprotected
baselines}, because a single linear map can fit all tasks jointly whenever
no key maps to two targets and capacity suffices. Forgetting must be
\emph{forced}: our tasks share a vocabulary subset whose successors differ
per task (interference bench), or belong to \emph{families} with correlated
symbols and partially shared transitions (long-sequence bench; 8 families
$\times$ 3 variants $=$ 24 tasks), optionally with family residuals
confined to a rank-3 subspace (low-rank bench). Retention is measured by
\emph{frozen-circuit probes}: a deep copy of the system with all
plasticity, fast-weight capture, and identity updates disabled, evaluated
on a past task after the full schedule. Forward transfer is measured as
time-to-criterion on a family's later variants relative to its first.

\section{Results}
\label{sec:results}

\subsection{Layered defenses on the interference bench}
Stripping mechanisms one at a time (600 ticks per phase, schedule
A$\to$B$\to$C$\to$A) yields post-transient recall error on task A of
\textbf{0.18} (full), 0.34 (no routing), 0.32 (no routing or cascade), and
\textbf{0.82} (dense unprotected baseline); relearned error returns to the
originally-learned level (0.042 vs.\ 0.045) for the full system. A
frozen-circuit probe of an earlier configuration showed task A's circuit
degrading only $0.085 \to 0.132$ across 1200 ticks of other tasks:
graceful power-law sag, not erasure.

\subsection{Long-sequence bench: retention, capacity, transfer}
\label{sec:longbench}

Table~\ref{tab:longbench} shows the full-rank bench (24 tasks, 500
ticks/task). Capacity grew to 8 columns for 24 tasks across all
conditions---tracking the 8 families, i.e.\ \emph{diversity}, not task
count. Retention is flat with age (oldest vs.\ newest third). Later family
variants reach criterion up to $2.6\times$ faster than first variants.

\begin{table}[ht]
\centering\small
\begin{tabular}{lrrrrrr}
\toprule
Encoder & TTC v1 & TTC v2+ & Fwd & Ret.\ (old) & Ret.\ (new) & Cols\\
\midrule
Fixed random               & 197 & 87  & $2.3\times$ & \textbf{0.214} & 0.216 & 8\\
Fixed $+$ MP               & 298 & 159 & $1.9\times$ & 0.253 & 0.255 & 8\\
Learnable, naive           & 240 & 129 & $1.9\times$ & 0.694 & 0.434 & 8\\
Learnable, consolidated    & 220 & 96  & $2.3\times$ & 0.217 & 0.224 & 8\\
\textbf{Consolidated $+$ MP} & 232 & \textbf{88} & $\mathbf{2.6\times}$ & 0.216 & \textbf{0.209} & 8\\
Consolidated $+$ MP $+$ decay & 220 & 89 & $2.5\times$ & 0.229 & 0.219 & 8\\
\bottomrule
\end{tabular}
\caption{Full-rank long-sequence bench (24 tasks; TTC = ticks to criterion;
Ret.\ = frozen-circuit error, oldest/newest third of tasks). The learned
dictionary with matching pursuit is the best overall system; the
$\mathrm{fixed}+\mathrm{MP}$ control confirms the mechanism (explaining
away pays only when atoms are learned to be worth explaining with).}
\label{tab:longbench}
\end{table}

\subsection{The representation-learning arc}
\label{sec:encoderarc}

Enabling encoder learning naively reintroduced catastrophic aging
(old-task retention error 0.69 vs.\ 0.21 fixed). Per-unit consolidation
with novelty recruitment rescued retention (0.29) but \emph{consistently
destroyed forward transfer} ($0.9\times$) across three allocation variants
(incremental, one-shot imprinting, residual imprinting). Instrumentation
localized the cause: learned features \emph{inherit input correlations}
that random projections destroy---within-task code separation degraded from
$|\!\cos\!| = 0.21$ (random) to $0.33$--$0.45$ (learned)---and the textbook
remedy, continuous anti-Hebbian decorrelation, would perpetually move
mature features, i.e.\ exactly the churn consolidation exists to prevent.
\textbf{The stability--plasticity dilemma re-emerges one level down, at the
representation layer, and consolidation alone does not solve it there.}

What solved it was architectural: the two-pathway residual code. With the
shared component absorbed by column expectations, the dictionary sees only
near-orthogonal residuals; the same consolidated learner then reached
parity with the fixed encoder on all task metrics with \emph{better} code
geometry ($0.036$ vs.\ $0.054$), and, after acquisition was decoupled from
allocation (atomic re-encoding plus maturity gating), surpassed it
(Table~\ref{tab:longbench}). The whole system also improved: acquisition
time halved ($427 \to 197$) and transfer rose from $1.6\times$ to
$2.3\times$ even for the fixed encoder.

\subsection{Low-rank structure, capacity limits, and expressivity}

Confining family residuals to a rank-3 subspace makes symbol directions
inherently correlated. Matching-pursuit inference plus the learned
dictionary beat the random encoder where theory requires---single-task
error floor 0.168 vs.\ 0.188; code separation 0.048 vs.\ 0.137---until
\textbf{capacity exhaustion}: with 288 symbols against 128 units, all
virgin rows were consumed within two of eight families, after which floors
degraded ($0.168 \to 0.258$) while the random encoder stayed flat (nothing
to exhaust; mediocre everywhere, forever). Usage decay converted the cliff
into a graceful trade-off (transfer $0.9\times \to 1.2\times$ at fixed
capacity). The residual transfer gap versus the fixed encoder was then
closed by the conjunctive expansion: on rank-3, learned$+$MP improved from
$1.2\times$ to $1.7\times$ forward transfer (TTC on later variants
$381 \to 155$), matching the fixed encoder---confirming the gap had been
readout expressivity, not representation quality.

\subsection{Metabolic embedding}
\label{sec:thermo}

\strata\ updates are outer products, so the Frobenius norm of every weight
change is available at no cost; the host organism charges this synaptic
work against a conserved energy economy in which work generates \emph{load}
(heat), load raises entropy, entropy degrades fidelity, and fidelity gates
the next tick's plasticity. Two results: (i) at a standard energy capacity,
baseline cognitive work alone pinned load above the learning-reachability
ceiling---the organ entered a saturated no-learning attractor (accuracy
0.12; learning unreachable on 397/400 ticks), reproducing a known
pathology of the host system in a learning organ; (ii) with a
brain-sized energy budget, a \textbf{burst/cool cadence emerged from pure
physics}: learning-reachability flickered at $\sim$4-tick spacing (86
ceiling crossings in 400 ticks) with no cooling rule anywhere in the code,
at an accuracy cost relative to the unconstrained control (0.34 vs.\ 0.57)
attributable to a separately documented saturated-fidelity attractor of the
host. Methodologically, the energy instrumentation acted as an autonomic
bug detector: a silent learning collapse (a missing normalization) was
caught because the organ went \emph{cold}.

\subsection{Scale verification at 48 tasks}
\label{sec:scale48}

Doubling the schedule (16 families $\times$ 3 variants, rank-3 residuals,
400 ticks/task) with the growing dictionary (initial size 64) and adaptive
claim bands produced the falsification cycle recorded in
Table~\ref{tab:falsification}: the first run spawn-stormed (48 columns for
48 tasks; newest-third frozen-probe error above 1.0 because tasks arriving
after the column cap were orphaned into no-learning fallback), the cause
was traced to background statistics absorbing transition swings, and after
the statistics-freeze fix both a static-threshold control and the adaptive
retest produced \textbf{exactly 16 columns for 16 families in every
condition}---diversity scaling intact at twice the task count---with
retention flat with age (best row: 0.306 oldest vs.\ 0.310 newest third)
and forward transfer of 1.4--1.6$\times$. Adaptive bands slightly
outperformed static thresholds after the fix (e.g., fixed encoder: 0.306
vs.\ 0.334 old-task retention). One honest finding, settled by a
pre-committed reading rule (any ambiguous outcome at under-demand capacity
mandates a rerun clearing demand before conclusions): at dictionary
capacities of 128, 512, and 768 units---the last comfortably clearing the
576-symbol allocation demand---the fixed encoder remains the best overall
row on rank-3 at 48 tasks. The gap is therefore \textbf{capacity-%
independent}: the learned-dictionary advantage of
Table~\ref{tab:longbench} is established on isotropic residuals and does
not extend to low-rank residual structure at this scale. Localizing why is
open work.

\subsection{Depth: a two-layer stack}
\label{sec:depth}

Layer 2 ticks only at Layer 1's regime boundaries: each detected
transition, after settlement, is pooled into one event
$u = \sum_t \nu_t\,\phi(z_t, e_t) / \sum_t \nu_t$, and L2 runs the identical
mechanism suite over events. Its predictions decode into bounded routing
priors and context seeding at the next transition. Across three seeds the
revisit transient (error in the first 40 ticks of a returning regime)
dropped \textbf{15--31\%} (mean $\approx 23\%$). Three lessons: (i)
\emph{emit what arrived, not what confusion felt like}---pooling the
surprise burst itself made all boundaries look identical and forbade L2
differentiation (fixing this raised the payoff from 18\% to 31\%); (ii)
absorption recurred at L2 as context--centroid co-drift over a shared event
stream (affinity 0.92 across a complete schedule change), with the
additional geometric fact that \emph{mixtures of unit vectors crowd
together}, making mixture-level identity intrinsically thin-margined; and
(iii) a measured \textbf{design constraint for hierarchy} that we propose
as a general law: \emph{hierarchical abstraction density requires geometric
timescale separation}. L2 \emph{can} differentiate regime-arrival types
(2--3 columns tracking schedule structure under event-granularity
thresholds), but at tens of events per run, splitting experience across
columns starves each successor map and prior quality degrades ($31\% \to
10\%$). The system thereby identifies, correctly, that the stream does not
yet contain enough regime-level experience to support differentiated
regime-level structure---the best overseer at small event budgets is the
least differentiated one---and quantifies how much more it would need.
Abstraction at level $N{+}1$ requires orders of magnitude more time than
level $N$; schedule-level structure belongs to a third layer, not to a
threshold adjustment at the second.

\section{The Falsification Record}
\label{sec:falsification}

Each mechanism was tested against a failure criterion stated \emph{before}
the experiment ran. We report the record in full, including the entries
where the design as first built failed its own test and had to be
recalibrated---these entries are, in our view, the strongest evidence that
the taxonomy is doing real work rather than post-hoc justification.

\begin{table}[ht]
\centering\small
\begin{tabular}{p{2.9cm} p{4.4cm} p{3.2cm} p{3.8cm}}
\toprule
\textbf{Mechanism} & \textbf{Pre-registered hypothesis} &
\textbf{Falsifier} & \textbf{Outcome}\\
\midrule
Naive encoder learning & Unprotected representation learning will
reintroduce forgetting (stated before implementation) & Retention remains
flat & \emph{Hypothesis confirmed}: old-task error $0.23 \to 0.76$\\
\addlinespace
Consolidation alone & Per-unit consolidation suffices for a learnable
encoder & Transfer or retention degrades & \textbf{Falsified} three times
(three allocation variants): retention rescued, transfer destroyed
($0.9\times$); cause measured as correlation inheritance\\
\addlinespace
Two-pathway residual coding & Removing shared structure before the
dictionary resolves correlation inheritance & Code separation fails to
recover & Survived: $|\!\cos\!|$ $0.33$--$0.45 \to 0.036$, better than
random ($0.054$); all task metrics improved\\
\addlinespace
Conjunctive expansion & The rank-3 transfer gap is readout expressivity,
not representation quality & Gap fails to close with superior codes &
Survived decisively: $1.2\times \to 1.7\times$, TTC $381 \to 155$\\
\addlinespace
Structural growth & Zero-initialized growth is backward-transparent &
Any existing code or readout changes & \textbf{Falsified as designed}
(fresh random rows won inference for old inputs); redesigned with dormant
capacity; then survived bit-exactly (unit-tested)\\
\addlinespace
Two-layer stack & Top-down priors cut the revisit transient & No
measurable reduction & Survived (15--31\%), after one interface
recalibration: pooling the surprise burst fed L2 undiscriminative events
(18\%); pooling post-settlement arrivals fixed it (31\%)\\
\addlinespace
L2 diversity scaling & L2 columns track schedules as L1 columns track
tasks & Column count fails to track schedule diversity & \textbf{Falsified
at current event budgets}; refined into the geometric-timescale design
constraint (\S\ref{sec:depth})\\
\addlinespace
Adaptive bands at scale & Statistics-tracked thresholds preserve diversity
scaling at 48 tasks / 16 families & Spawn storm or threshold collapse &
\textbf{Falsified as designed}: spawn storm (48 columns for 48 tasks;
family sharing collapsed; late-task retention pathological). Cause traced
to background statistics updating during transition swings, inflating the
adaptive floor. Fix: statistics freeze during transitions (the same law as
centroids and claim statistics). \emph{Revalidation passed}: 16 columns for
16 families, matching a static-threshold control run
(\S\ref{sec:scale48})\\
\addlinespace
External adjudication, stage 1 & The trichotomy decomposes forgetting on a
benchmark we did not design (Split-MNIST, class-incremental; naive
fine-tuning, EWC, SI, replay) & Any event passing all three exclusions, or
$>$30\% of forgetting magnitude unattributed & Survived: 33/33 events
attributed across two seeds; substrate recovery $=1.00$ on \emph{every}
event (all forgetting parameter-carried in these baselines); basis
alignment 0.65--0.92 against null floors $\le$0.37; timescale attributed
on only 2/33 (at full data, $10\times$-slowed training still converges to
the interfering solution). Disclosed per Amendment 2: every event had at
least one vacuous layer, partially jamming the basis door to (c) on this
architecture; the verdict rests chiefly on the fully-testable substrate
exclusion. Prediction (b)'s candidate modes were structurally untestable
here (plain-SGD MLPs carry no normalization or optimizer state); that is
the CIFAR stage's question\\
\bottomrule
\end{tabular}
\caption{The falsification record. Bold outcomes are failures of the
design as first built; every failure produced either a redesign that then
passed, or a quantified constraint.}
\label{tab:falsification}
\end{table}

\section{Related Work}

Regularization approaches estimate parameter importance and slow important
weights \cite{kirkpatrick2017,zenke2017}; \strata's cascade and per-unit
consolidation play a similar role but are computed online from local
signals, with no Fisher estimation and no task boundaries, and are
demonstrated to be insufficient alone at the representation layer
(\S\ref{sec:encoderarc}). Replay methods \cite{shin2017} revisit stored or
generated data; \strata's calm-time distillation replays only its own
fast-weight content. Structural methods \cite{rusu2016} isolate tasks in
grown modules; \strata's columns are similar in spirit but are recruited by
measured orphanhood rather than task labels, grow with diversity rather
than task count, and share a learned encoder. The fast/slow organization
follows complementary learning systems \cite{mcclelland1995}; the cascade
follows \cite{fusi2005,benna2016}; sparse expansion coding follows
\cite{olshausen1996,dasgupta2017}; residual coding follows predictive
coding \cite{rao1999}; matching pursuit is classical \cite{mallat1993};
fixed high-dimensional dynamics with learned readouts follow reservoir
computing \cite{maass2002,jaeger2001}. Our contribution is the
composition under strict continual-learning constraints, and the failure
taxonomy that the composition exposed.

\section{Limitations}
\label{sec:limitations}

We state these plainly. (1)~\textbf{Scale}: all results are at $d{=}24$,
$\le$48 tasks, $\sim$300 symbols, linear readouts over fixed expansions;
per-tick dynamics are interpreted, not optimized. (2)~\textbf{Benchmark
ownership}: the \strata\ benches were designed by us. The pre-registered
external protocol (\code{paper/preregistration.md}) has now completed its
first stage---Split-MNIST with standard baselines, reported in the
falsification record---with the trichotomy surviving; the harness and
hyperparameters were still ours, the architectures tested carry no
normalization or optimizer state, and the Split-CIFAR stage with richer
architectures remains open.
(3)~\textbf{Depth}: one representation layer plus a two-layer routing
hierarchy; whether the stability machinery composes through deep feature
hierarchies is untested and is, in our judgment, the single most important
open question. (4)~\textbf{Hand-set constants}: routing thresholds were
tuned to measured geometry; the adaptive-band mechanism \emph{failed its
first pre-registered scale falsifier} (spawn storm;
Table~\ref{tab:falsification}) and passed revalidation only after the
statistics-freeze redesign (\S\ref{sec:scale48}); further scale remains
untested. (5)~\textbf{Capability
ceiling}: stability was purchased with deliberately weak learners; the
learned-versus-fixed dictionary advantage, while now positive, is modest.
(6)~Some forgetting is information-theoretically mandatory at finite
capacity on unbounded streams; \strata\ chooses \emph{which} (graceful,
least-used-first), not \emph{whether}.

\section{Conclusion}

Treated as an architectural default rather than a training pathology,
catastrophic forgetting decomposes into engineering constraints that can be
named, measured, unit-tested, and removed. Seven local mechanisms suffice,
at prototype scale, for label-free, replay-free continual learning with
flat retention, diversity-scaled capacity, positive forward transfer,
graceful forgetting, metabolic self-regulation, and---after resolving the
representation-layer dilemma architecturally---representation learning that
helps rather than harms. The failure taxonomy generalized upward through
one level of hierarchy, in level-specific costume, exactly as a useful
theory should: it told us where to look when things broke. The decisive
next experiments are composition under depth, benchmarks we do not own,
and an order of magnitude of scale.

\paragraph{Reproducibility.} All experiments are reproducible from the
accompanying code: \code{python -m strata.demo} (interference bench),
\code{python -m strata.longbench} (long-sequence bench; \code{--rank 3},
\code{--conjunctive 512}, \code{--grow-from}, \code{--adaptive}),
\code{python -m strata.stack} (depth bench), \code{python main.py --cortex
[--thermo]} (metabolic embedding), and \code{pytest} (62 tests, including
exact backward-transparency of structural growth and the zero-plasticity
freeze contract).

\begin{thebibliography}{99}
\bibitem{benna2016} M.~K. Benna and S.~Fusi. Computational principles of
synaptic memory consolidation. \emph{Nature Neuroscience}, 19:1697--1706,
2016.
\bibitem{dasgupta2017} S.~Dasgupta, C.~F. Stevens, and S.~Navlakha. A
neural algorithm for a fundamental computing problem. \emph{Science},
358(6364):793--796, 2017.
\bibitem{french1999} R.~M. French. Catastrophic forgetting in connectionist
networks. \emph{Trends in Cognitive Sciences}, 3(4):128--135, 1999.
\bibitem{fusi2005} S.~Fusi, P.~J. Drew, and L.~F. Abbott. Cascade models of
synaptically stored memories. \emph{Neuron}, 45(4):599--611, 2005.
\bibitem{jaeger2001} H.~Jaeger. The ``echo state'' approach to analysing
and training recurrent neural networks. GMD Technical Report 148, 2001.
\bibitem{kirkpatrick2017} J.~Kirkpatrick et al. Overcoming catastrophic
forgetting in neural networks. \emph{PNAS}, 114(13):3521--3526, 2017.
\bibitem{maass2002} W.~Maass, T.~Natschl\"ager, and H.~Markram. Real-time
computing without stable states. \emph{Neural Computation},
14(11):2531--2560, 2002.
\bibitem{mallat1993} S.~Mallat and Z.~Zhang. Matching pursuits with
time-frequency dictionaries. \emph{IEEE Transactions on Signal Processing},
41(12):3397--3415, 1993.
\bibitem{mcclelland1995} J.~L. McClelland, B.~L. McNaughton, and R.~C.
O'Reilly. Why there are complementary learning systems in the hippocampus
and neocortex. \emph{Psychological Review}, 102(3):419--457, 1995.
\bibitem{mccloskey1989} M.~McCloskey and N.~J. Cohen. Catastrophic
interference in connectionist networks. \emph{Psychology of Learning and
Motivation}, 24:109--165, 1989.
\bibitem{olshausen1996} B.~A. Olshausen and D.~J. Field. Emergence of
simple-cell receptive field properties by learning a sparse code for
natural images. \emph{Nature}, 381:607--609, 1996.
\bibitem{rao1999} R.~P.~N. Rao and D.~H. Ballard. Predictive coding in the
visual cortex. \emph{Nature Neuroscience}, 2(1):79--87, 1999.
\bibitem{rusu2016} A.~A. Rusu et al. Progressive neural networks.
\emph{arXiv:1606.04671}, 2016.
\bibitem{shin2017} H.~Shin, J.~K. Lee, J.~Kim, and J.~Kim. Continual
learning with deep generative replay. \emph{NeurIPS}, 2017.
\bibitem{zenke2017} F.~Zenke, B.~Poole, and S.~Ganguli. Continual learning
through synaptic intelligence. \emph{ICML}, 2017.
\end{thebibliography}

\appendix

\section{Hyperparameters}
Default configuration (dimensionless units; one sample per tick):
input $d{=}24$; code $D{=}128$, sparsity $s{=}8$; state $n{=}32$ per
column; $\rho_{\mathrm{context}}{=}0.08$;
$\theta_{\mathrm{route}}{=}0.65$, $\theta_{\mathrm{claim}}{=}0.8$;
spawn: leaky rate 0.15, threshold 0.8; cascade: 4 levels,
$\tau_1{=}8$, ratio 4; learning rates: state readout 0.05, code readout
0.3, distillation 0.1, fast weights 0.6 (novelty-gated, hard floor at 5\%
of maximum rate); CUSUM slack 0.5, threshold 8, transition refractory 15
ticks, mature-column damping 0.15; encoder learning 0.15 with
consolidation $\kappa{=}0.2$, recruitment threshold 0.5, usage decay
$2\times10^{-4}$ (when enabled); conjunctive expansion 512 with $S{=}12$
(when enabled); layer-2: event-driven, settle window 15 ticks,
$\rho{=}0.12$, $\theta_{\mathrm{claim}}{=}0.9$, cascade $\tau_1$ offset
$\times16$.

\section{On the use of failure-driven development}
Ten of the mechanisms reported here were not in the original design; each
was added after an instrumented experiment falsified the design as it
stood, and three (absorption hysteresis, dormant growth, post-settlement
event pooling) were forced by failures that the a-priori mathematics gave
no reason to expect. We record this because it bears on methodology:
in this domain, contact with data was worth more than derivation at every
single decision point.

\end{document}
